\subsection{buffer\_pool}\label{buffer_pool}
The \verb|buffer_pool| module manages a fixed number of pre-allocated
\verb|buffer_t| (Section \ref{buffer}) and hands them out on a used/free
basis. Callers obtain a free buffer with \verb|buffer_pool_get()| and return
it after use with \verb|buffer_pool_return()|, instead of repeatedly
allocating and freeing individual buffers. This avoids heap fragmentation and
makes the memory footprint deterministic at run time.

\subsubsection*{State}
The pool is described by \verb|buffer_pool_t| (Listing \ref{buffer_pool_t}).
\verb|buffer| is the array of the managed \verb|buffer_cnt| buffers;
\verb|type| denotes the operating mode of the pool and \verb|next| the index
of the next element (negative if a linear buffer is used in a ring).
\verb|sbuffer| is a second field, currently reserved for extensions.

\begin{listing}
\begin{minted}{C}
typedef struct buffer_pool_s {
    uint8_t buffer_cnt;  // number of buffers in the pool
    buffer_t **buffer;   // array of buffers
    buffer_t **sbuffer;  // reserved
    b_type_e type;       // operating mode of the pool
    int8_t next;         // index of the next element
} buffer_pool_t;
\end{minted}
\caption{Definition of the \texttt{buffer\_pool\_t} structure}
\label{buffer_pool_t}
\end{listing}

\subsubsection*{API}
{\sloppy\emergencystretch=3em
\begin{description}
\item[\texttt{buffer\_pool\_new(lcnt, charCnt, line\_type, type)}] Creates a
  pool of \verb|lcnt| buffers of \verb|charCnt| bytes each. \verb|line_type|
  sets the type of the individual buffers (\verb|LINEAR| or \verb|RING|),
  \verb|type| the operating mode of the pool. If a partial allocation fails,
  already created buffers are freed again and \verb|NULL| is returned.
\item[\texttt{buffer\_pool\_free(pool)}] Frees all buffers and the pool itself.
\item[\texttt{buffer\_pool\_get(bp)}] Returns the first free buffer, marks it
  \verb|BUFFER_USED| and returns it; \verb|NULL| if no buffer is free.
\item[\texttt{buffer\_pool\_return(bp, buffer)}] Returns a buffer to the pool.
  The index range, the membership of the pool (pointer comparison) and that
  the buffer is no longer in use are checked; it is then reset with
  \verb|buffer_reset()|. A buffer still marked \verb|BUFFER_USED| is rejected
  with \verb|EM_ERR| -- so the content must have been consumed beforehand
  (e.g.\ via \verb|buffer_get()|).
\item[\texttt{buffer\_pool\_print(bp, title)}] Prints the pool's buffers for
  debugging.
\end{description}
\par}
